% ============================================================================
% Lab Chapter for Part VII — Practicals 1–6
% ============================================================================
\chapter{Part VII --- Hands-On Practicals}
\label{ch:lab7}

\bytetip[fun]{Let's go online --- safely! These practicals teach you to browse, search, email, and stay protected while exploring the digital world.}

\begin{labactivity}{1}{Exploring a Web Browser}{10 min}
\youwillneed{\faDesktop\ Computer with Internet \quad \faClock\ 10 min \quad \faChrome\ Browser}
\textbf{Objective:} Navigate a web browser and identify its parts.\\[4pt]
\textbf{Steps:}
\begin{enumerate}[label=\protect\stepcircle{\arabic*}, leftmargin=2cm, itemsep=4pt]
  \item Open Google Chrome (or your school's default browser).
  \item Identify: Address Bar, Tabs, Refresh Button, Back/Forward Buttons.
  \item Type \texttt{www.wikipedia.org} in the address bar and press Enter.
  \item Click on any link to navigate to a new page. Use the Back button.
  \item Open a new tab (\key{Ctrl}+\key{T}) and visit \texttt{www.google.com}.
\end{enumerate}
\textbf{Expected Outcome:} Students navigate browser elements confidently.
\welldone
\end{labactivity}

\begin{labactivity}{2}{Smart Searching}{15 min}
\youwillneed{\faDesktop\ Computer with Internet \quad \faClock\ 15 min \quad \faSearch\ Google}
\textbf{Objective:} Use advanced search techniques.\\[4pt]
\textbf{Steps:}
\begin{enumerate}[label=\protect\stepcircle{\arabic*}, leftmargin=2cm, itemsep=4pt]
  \item Search: ``planets in the solar system'' --- note the results.
  \item Search with quotes: \texttt{"largest planet in the solar system"} --- compare results.
  \item Search with \texttt{site:}: \texttt{site:wikipedia.org Indian festivals}.
  \item Search with minus: \texttt{python -snake} (to find the programming language).
  \item For each search, write down how many results appeared and the top result.
\end{enumerate}
\textbf{Expected Outcome:} Students use 4 search techniques effectively.
\welldone
\end{labactivity}

\begin{labactivity}{3}{Bookmarking Websites}{10 min}
\youwillneed{\faDesktop\ Computer with Internet \quad \faClock\ 10 min \quad \faChrome\ Browser}
\textbf{Objective:} Save and organise favourite websites.\\[4pt]
\textbf{Steps:}
\begin{enumerate}[label=\protect\stepcircle{\arabic*}, leftmargin=2cm, itemsep=4pt]
  \item Visit 3 educational websites (e.g., Wikipedia, National Geographic Kids, Khan Academy).
  \item Bookmark each one: click the star icon in the address bar.
  \item Create a bookmark folder called ``School Resources.''
  \item Move all 3 bookmarks into that folder.
  \item Access a bookmarked site from the bookmarks bar.
\end{enumerate}
\textbf{Expected Outcome:} Students save and organise 3 bookmarks.
\welldone
\end{labactivity}

\begin{labactivity}{4}{Composing an Email}{15 min}
\youwillneed{\faDesktop\ Computer with Internet \quad \faClock\ 15 min \quad \faEnvelope\ Email account}
\textbf{Objective:} Write and send a properly formatted email.\\[4pt]
\textbf{Steps:}
\begin{enumerate}[label=\protect\stepcircle{\arabic*}, leftmargin=2cm, itemsep=4pt]
  \item Open your email (Gmail, Outlook, or school email).
  \item Click ``Compose'' (or ``New Email'').
  \item In the ``To'' field, type your teacher's email address.
  \item Subject: ``Homework Query --- [Your Name].''
  \item Body: Write a polite message asking about a homework deadline.
  \item Add your name as a sign-off. Click Send (or save as Draft).
\end{enumerate}
\textbf{Expected Outcome:} Students compose a well-structured email.
\welldone
\end{labactivity}

\begin{labactivity}{5}{Identifying Phishing Emails}{10 min}
\youwillneed{\faDesktop\ Computer \quad \faClock\ 10 min \quad \faFile\ Worksheet}
\textbf{Objective:} Spot the signs of phishing emails.\\[4pt]
\textbf{Steps:}
\begin{enumerate}[label=\protect\stepcircle{\arabic*}, leftmargin=2cm, itemsep=4pt]
  \item Your teacher will show you 5 sample emails (printed or on screen).
  \item For each email, decide: REAL or PHISHING?
  \item Write down the clues that helped you decide.
  \item Discuss your answers with the class.
  \item Create a ``Phishing Red Flags'' checklist with at least 5 warning signs.
\end{enumerate}
\textbf{Expected Outcome:} Students correctly identify 4 out of 5 phishing emails.
\welldone
\end{labactivity}

\begin{labactivity}{6}{Digital Citizenship Pledge}{15 min}
\youwillneed{\faDesktop\ Computer \quad \faClock\ 15 min \quad \faFileWord\ MS Word or Paper}
\textbf{Objective:} Create a personal digital citizenship pledge.\\[4pt]
\textbf{Steps:}
\begin{enumerate}[label=\protect\stepcircle{\arabic*}, leftmargin=2cm, itemsep=4pt]
  \item Open MS Word (or use paper).
  \item Write a title: ``My Digital Citizenship Pledge.''
  \item Write 8 promises about safe, kind, and responsible internet use.
  \item Add your name, date, and signature at the bottom.
  \item Decorate it and print it (or display it on screen).
  \item Share your pledge with the class.
\end{enumerate}
\textbf{Expected Outcome:} A personalised digital citizenship pledge document.
\welldone
\end{labactivity}

\vspace{12pt}
\begin{center}
  \qrcode[height=2cm]{https://example.com/module7}\\[6pt]
  {\small\sffamily\textcolor{NavyBlue}{Scan for extra Part VII resources and activities!}}
\end{center}
